\section{Machine-learning timing correction}

The primary comparison evaluates the detector-shared antisymmetric dense MLP against the fixed paired CNN reference based on Onishi et al.~\cite{onishi2022}. A parallel locally connected MLP is also implemented to test a stronger temporal inductive bias without imposing convolutional weight sharing.

\begin{table}[H]
    \centering
    \caption{The two machine-learning models retained in the final waveform pipeline.}
    \label{tab:model-summary}
    \small
    \begin{tabular}{llll}
        \toprule
        Model & Input representation & Prediction structure & Scientific role \\
        \midrule
        Antisymmetric MLP & \(s_1,s_2\) & \(g_\theta(s_1)-g_\theta(s_2)\) & Proposed detector-shared model \\
        Onishi CNN & stacked \([2,\mathrm{time}]\) & \(f_\theta(s_1,s_2)\) & Fixed paired-CNN reference \\
        \bottomrule
    \end{tabular}
\end{table}


\subsection{Detector-shared antisymmetric MLP}\label{sec:mlp-model}

The proposed model applies the same dense detector scorer \(g_\theta\) independently to the two aligned waveforms:
\begin{equation}\label{eq:mlp-correction}
    y_\theta(s_1,s_2)
    =
    g_\theta(s_1)-g_\theta(s_2).
\end{equation}
The two branches share all parameters. Exchanging the detector inputs therefore gives
\begin{equation}
    y_\theta(s_2,s_1)
    =
    -y_\theta(s_1,s_2),
\end{equation}
so detector-swap antisymmetry is enforced exactly by construction.

In the current model space, \(g_\theta\) is a fully connected network with hidden widths \(128\) and \(64\), SiLU activations, and one scalar output. The candidate grid varies the learning rate between \(10^{-3}\) and \(3\times10^{-4}\) with batch size 128. Training uses stochastic gradient descent with Nesterov momentum (momentum coefficient 0.9), RMSE loss, up to 600 epochs, RMSE-based early stopping on an internal holdout, gradient clipping, and an additional low-gradient stopping condition. When several hyperparameter candidates are available, validation CTR selects the configuration and a fresh final model is then trained on the full development population using the same internal early-stopping mechanism.

A dense scorer is used intentionally. Once each waveform has been aligned to the LED crossing, absolute temporal position within the retained window is physically meaningful, so translation equivariance is not imposed as a prior on the proposed detector-level model.

\reportfigure
    {figures/mlp-architecture.pdf}
    {Proposed antisymmetric MLP. The same dense scorer is applied independently to \(s_1\) and \(s_2\), and the two scalar detector scores are subtracted.}
    {fig:mlp-architecture}

\subsection{Locally connected antisymmetric MLP}
The locally connected variant preserves the same paired prediction
\[
    y_{\theta}(s_1,s_2)=g_{\theta}(s_1)-g_{\theta}(s_2),
\]
but restricts the first layer to contiguous temporal receptive fields. Each receptive field produces one scalar node with its own weight vector and bias. Consecutive fields may overlap, with stride equal to the receptive-field width minus the configured overlap. The weights are position specific rather than shared, distinguishing the layer from a convolution while still encoding the prior that early representations should depend only on local waveform structure. The local outputs are followed by an ordinary dense stack. The complete configured temporal grid is retained for this model so that receptive fields always contain consecutive native samples.

\subsection{Onishi paired CNN reference}\label{sec:onishi-model}

The reference model follows the paired-waveform CNN used by Onishi et al.~\cite{onishi2022}. It receives the detector pair as a stacked \([2,\mathrm{time}]\) input and predicts one pair-level correction,
\begin{equation}\label{eq:onishi-correction}
    y_\theta=f_\theta(s_1,s_2).
\end{equation}

The implemented reference architecture is fixed. It uses a \(2\times5\) Conv2D layer with 32 channels and ReLU, followed by two \(1\times3\) Conv2D layers with 32 and 64 channels, respectively, also with ReLU. The resulting features are flattened and passed through a dense layer with 256 units and ReLU before the scalar output. The first convolution spans both detector rows and therefore fuses the detector pair immediately. Unlike Eq.~\eqref{eq:mlp-correction}, detector-swap antisymmetry is not enforced by the architecture.

The current reference training configuration uses Adam, MSE loss, batch size 128, an initial learning rate of \(10^{-3}\), and 600 epochs. The learning rate is reduced by a factor of ten at epochs 180 and 360. Because this is a fixed reference configuration, no architecture or optimizer hyperparameter is selected from the validation CTR.

\subsection{Controlled comparison across waveform windows}

The models are executed as independent \texttt{model\_study} runs. Each run evaluates the same configured waveform windows and persists the blind residuals, model outputs, split identities, XAI arrays, and result tables.

The completed model-study runs are then combined with the repository's \texttt{compare-runs} workflow. Before comparison, the workflow checks compatible window definitions and voltage sets. Model-output correlation is computed only after aligning the persisted blind-test event identities.

This separation is useful operationally because the dense MLP, locally connected MLP, and Onishi CNN can be trained independently, while the final comparison still operates on compatible persisted blind-test results.

\subsection{Interpretability and model-output agreement}

All three implemented models expose the same gradient-based temporal sensitivity measure: the absolute input gradient is averaged over the evaluated events and detector rows and projected onto the retained waveform time axis~\cite{simonyan2014}. The resulting profile identifies where a fitted model is most sensitive, but it is not interpreted by itself as proof of causal timing information.

The run comparison additionally computes pairwise Pearson correlations between model blind-test outputs after event-ID alignment. This quantity answers a complementary question: whether two models with different representations are learning similar event-by-event corrections even when their final CTR values are close.
